\scanpage{87}
% archive sheet number: 86
En prenant $G={\unsure{G|\mathfrak S|\mathcal S}}_3$, $\gamma=\widehat{\mathbb Z}^{\unsure{\times|*}}$ comme dessus, on trouve donc
\diagram{\begin{tikzcd}[ampersand replacement=\&,column sep=large,row sep=large]
G_3*\gamma \arrow[d, "{\substack{\text{via l'op. de }G_3\text{ sur }\unsure{\Pi|\Gamma|\Omega}\\ \text{et }\rho_{\varepsilon1.}}}" description] \&[-1em] \simeq \&[-1em] G_3\underset{1/2}{\cdot}\gamma^{(*G_3)}\\
\Omega \& \&
\end{tikzcd}}
% slanted
\margin{On peut prendre ici les \unsure{gr.} au sens profini \ldots}

Cet hom n'est pas injectif -- on a des relations, \deleted{\unsure{qui|pour}} \inserted{notamment l'une} exprimant la commutation des $\rho_{\varepsilon1.}$ et $\rho_{1\varepsilon'.}=\pi\rho_{\varepsilon.1}\pi^{-1}$ ($\pi=\pi_{01}$ transposition de $0$ et $1$) donc on va définir dans $\gamma^{(*G_3)}$ un \deleted{sous}groupe quotient inv. par $G_3$, en divisant par les relations exprimant
\[
\rho^{\pi}_{\deleted{\unsure{\varepsilon}}}\varepsilon\ \text{commute à}\ \rho^1(\varepsilon')\qquad
(\varepsilon,\varepsilon'\in\gamma)
\]
On trouve
\[
G_3\underset{1/2}{\cdot}\,\overset{\overset{\simeq}{\longrightarrow}}{\unsure{E|\mathcal E}}
(\gamma^2*\gamma^2*\gamma^2)=(\gamma^2)^{(*G_3/\unsure{\pi|2}}
\]
où $E$ est le produit libre \inserted{(profini)} des trois copies de $\gamma^2$, permutées transitivt entre elles par les op. de $G_3$, et un hom
\diagram{\begin{tikzcd}[ampersand replacement=\&,column sep=small]
G_3\underset{1/2}{\cdot}\deleted{\unsure{\gamma}}(\gamma\times\gamma)^{(*G_3/\pi_{01})}
 \arrow[r]\arrow[dr] \& \Omega^+ \arrow[d,dash] \&[-1em] \subset \&[-1em] \Omega \\
\& \widetilde{\Omega}^{+} \&[-1em] \subset \&[-1em] \widetilde{\Omega}\arrow[u,phantom,"\rotatebox{90}{$\subset$}" description]
\end{tikzcd}}
Il serait intéressant d'examiner les propriétés d'inj. et surjectivité de ces flèches, de $=$ que les variantes discrètes
\margin{$\gamma_0=\{\pm1\}$}
\diagram{\begin{tikzcd}[ampersand replacement=\&,column sep=small]
G_3\underset{1/2}{\cdot}(\gamma_0\times\gamma_0)^{(*G_3/\pi_{01})}
 \arrow[r]\arrow[dr] \& \Omega_0^+\arrow[d] \&[-1em] \subset \&[-1em] \Omega_0\arrow[d] \\
\& \widetilde{\Omega}_0^+ \&[-1em] \subset \&[-1em] \widetilde{\Omega}_0
\end{tikzcd}}
